\section{Introduction}
\label{sec:intro}

Joint Embedding Predictive Architecture (JEPA) has been the subject of much active development. It
predicts in latent/embedding space rather than reconstructing pixels, discarding
irrelevant or unpredictable detail to focus on high-level structure. However, almost nothing is known rigorously about what its training dynamics do to the
underlying data structure.

We find that JEPA training dynamics themselves partition every feature direction by its true
predictive relationship to the target. We prove that in the linear JEPA
context, we get signed, per-feature regression, decodable from the
encoder's internal weight trajectory. We have an explicit signed
estimator computable from the trajectory alone, and an easy to deploy
algorithm we call \textsc{PlateauRecover} (code provided). Concretely:
the sign of $\rhostar$ is always recoverable from the trajectory
$\sigmar(t)$ alone, and on the positive branch so is its magnitude
(given only the eigenbasis, no further population-covariance input
needed) at rate $O(\eps^{1/L}|\log\eps|)$, a corollary of the
results in \S\ref{sec:estimating-spectrum}
(Theorem~\ref{thm:headline-formal}).

\label{ssec:positioning}
In the linear case, everything we recover here is something OLS (or
other gradient-descent methods) could already give us. We care about
it because it is a critical stepping stone to the non-linear regime of
representation learning on unlabeled data. There, there may be no
clean $y$, no fixed $\Sigma^{xx}$-type inner product, and no
normal-equations solver to fall back on. In particular, the
LeWorldModel architecture of \citet{Maes2026LeWM} is exactly such a
case. There, the training trajectory is the only handle on the
underlying regression structure: $\{\rhostar\}$ can be read off a
JEPA training trajectory without ever forming $\Sigma^{xx}, \Sigma^{yx}$.

\subsection*{Notation}

We fix notation for the covariances and the regression spectrum they define.
For data pairs $(x,y)\in\RR^d\times\RR^d$, write
$\Sigma^{xx} := \EE[xx^\top] \succ 0$ for the input covariance and
$\Sigma^{yx} := \EE[yx^\top]$ for the target-input cross-covariance.
The regression structure we recover throughout the paper is read off the
generalized eigenproblem
\[
\Sigma^{yx} v_r^* = \rho_r^* \Sigma^{xx} v_r^*.
\]
Unlike an ordinary eigenvector, $v_r^*$ is fixed only up to an
arbitrary rescaling. The generalized eigenvalue $\rho_r^*$ itself is
untouched by this rescaling, since it is a property of the
eigen-direction rather than of any particular representative $v_r^*$.
What does move is the $\Sigma^{xx}$-quadratic form
$\mu_r := v_r^*\cdot\Sigma^{xx}v_r^*$ of $v_r^*$ (guaranteed $>0$ since
$\Sigma^{xx}\succ0$), and with it the projected covariance
$\lambda_r^* := \rho_r^*\mu_r$. This invariance makes $\rho_r^*$ the
signed regression coefficient the rest of the paper tracks. The
collection $\{\rho_r^*\}_{r=1}^d$ (one signed scalar per feature
direction) is the ``regression spectrum'' of the pencil
$(\Sigma^{yx}, \Sigma^{xx})$ we are after. We package the triple
$(\rhostar, \mu_r, v_r^*)$ into a single labeled object that fixes this
scale-dependence once and for all:

\begin{definition}[Signed generalised eigenbasis]
\label{def:eigenbasis}
A \emph{signed generalised eigenbasis} for the pencil $(\Syx, \Sxx)$ is
a tuple $\{(\rhostar, \mu_r, v_r^*)\}_{r=1}^{d}$ with
$\rhostar \in \RR$, $v_r^* \in \RR^d \setminus \{0\}$,
$v_r^* \cdot \Sxx v_s^* = 0$ for $r \neq s$, and
\begin{equation}
\label{eq:gen-eigenpair}
\Syx v_r^* = \rhostar \Sxx v_r^*,
\qquad
\mu_r := v_r^* \cdot \Sxx v_r^*,
\qquad
\lamstar := \rhostar \mu_r.
\end{equation}
$\mu_r > 0$ is \emph{not} normalised to $1$. It is the free
$\Sxx$-quadratic form of the (arbitrarily scaled) eigenvector $v_r^*$,
matching \lean{Basic.SignedGenEigenpair.hmu_def}. Rescaling $v_r^*$
rescales $\mu_r$ and $\lamstar$ together while leaving
$\rhostar = \lamstar/\mu_r$ fixed, so $(\lamstar, \mu_r)$ is only
meaningful jointly, given a fixed scale for $v_r^*$. $\rhostar$ itself
is the scale-invariant quantity this paper recovers.
\end{definition}

\label{ssec:linear-jepa}
We now describe the model whose training we study: fix dimension $d$
and depth $L\ge2$ throughout (equal input/output dimension; the
rectangular case is left to future work). A depth-$L$ linear JEPA
factors its encoder as $\bar W = W_L\cdots W_1$,
each factor $W_a\in\RR^{d\times d}$ kept separate (never collapsed into
one matrix), plus a single unfactored predictor $V\in\RR^{d\times d}$.
This model trains to minimize the squared reconstruction error
\[
\mathcal{L}(\bar W, V) \;=\; \tfrac{1}{2}\,\EE\!\left[\big\|y - V \bar W x\big\|_2^2\right].
\]

This same eigenbasis is where the resulting training trajectory lives:
under aligned balanced initialisation, the encoder's weight matrix
$\Wbar(t)$ stays diagonal in it for all $t \ge 0$. Its diagonal entry
$\sigmar(t) := v_r^{*\top}\Wbar(t)v_r^*$ is the amplitude of feature
$r$ at time $t$. It obeys a Bernoulli-type ODE decoupled across
features at leading order, stated precisely as
Proposition~\ref{prop:diagonal-ode} in \S\ref{sec:diagonal-ode}
(equation~\eqref{eq:diagonal-ode}). Recovering $\{\rhostar\}$ means
understanding what training under this loss does to $\bar W$ and $V$
over time. This is where population gradient flow comes in.

\subsection*{Assumptions}
\label{ssec:assumptions-intro}

We work throughout under the following standing assumptions.

\begin{enumerate}[label=(A\arabic*),ref=A\arabic*,leftmargin=2.5em,itemsep=4pt]

\item \label{ass:population-flow}
\textbf{Exact population gradient flow.} Training proceeds under
\begin{equation}
\label{eq:jepa-flow}
\dot{\Wbar} = -\nabla_{\Wbar} \mathcal{L}
\ \ \text{(shorthand for the preconditioned flow of Lemma~\ref{lem:grad-flow})},
\qquad
\dot{V} = -\nabla_{V} \mathcal{L}.
\end{equation}
This is the theoretical idealisation; Algorithm~\ref{alg:plateau-recover}
(\S\ref{sec:algorithm}) instead runs discretised SGD on finite samples.

\item \label{ass:aligned-init}
\textbf{Aligned balanced initialisation.} We restrict throughout to
initialisation at scale $\eps\in(0,1)$:
\begin{equation}
\label{eq:init}
\bar W(0) = \eps \cdot U_x^\top U_x,
\qquad
V(0) = \eps \cdot U_y U_x^\top,
\end{equation}
where $U_x$ is the orthogonal basis of generalised eigenvectors of the
pencil $(\Sigma^{yx}, \Sigma^{xx})$ and $U_y$ its $\Sigma^{xx}$-orthonormal
image under $\Sigma^{yx}$ (Definition~\ref{def:eigenbasis}). The
balance condition is preserved under \eqref{eq:jepa-flow}
\citep{Saxe2014,Arora2018}. Exact alignment cannot hold in practice ---
the eigenbasis is unknown at initialisation; every identifiability result
here is stated under \eqref{eq:init} exactly, with rate degradation under
non-aligned initialisation not established.

\item \label{ass:gap}
\textbf{Strict spectral gap.} The signed values $\{\rhostar\}_{r=1}^d$
are strictly distinct and we order them as
$\rho_1^* > \rho_2^* > \cdots > \rho_d^*$. Write
$\delta_r := \min_{s \neq r}|\rhostar - \rho_s^*|$ for the gap and
$\delta := \min_r \delta_r$ for the minimum gap. Needed only for the
finite-sample chain (\S\ref{sec:finite-sample}); single-feature,
pure-trajectory results need no gap condition.

\item \label{ass:traj-reg}
\textbf{Trajectory regularity.} There exist constants $K_1, c_w > 0$,
independent of $\eps$, such that along the JEPA trajectory on
$[0, t_{\max}^*]$: (i) the encoder moves slowly,
$\|\dot{\Wbar}(t)\|_F \le K_1\,\eps^2$; (ii) the diagonal amplitudes of
Definition~\ref{def:eigenbasis} stay bounded below,
$\sigmar(t) \ge c_w\,\eps^{1/L}$. Both hold with equality up to constants
at $t=0$ under~\eqref{eq:init}. We do not re-derive their propagation to
every $t \in [0, t_{\max}^*]$ here; instead we take them as the standard
trajectory-regularity condition that informal ODE learning-dynamics
analyses typically leave implicit. This matches the hypothesis
structure of the companion's own \lean{bootstrap_consistency}
(\lean{hWbar_slow}, \lean{hdiag_t} in \texttt{BootstrapLemmas.lean}).

\end{enumerate}

\subsection*{Proof strategy}
\label{ssec:proof-strategy}

We analyze the global training dynamics by applying gradient flow at the factor
level---to each matrix $W_a$ individually and to $V$ directly---rather than to the
end-to-end product $\bar W$ as one matrix. This factorized parameterization induces
the depth-dependent preconditioned flow of \eqref{eq:jepa-flow}
(Lemma~\ref{lem:grad-flow}), yielding the characteristic Saxe-style depth exponents.
The proof reduces this coupled matrix system to scalar trajectory recovery in four
steps.

The predictor $V$ relaxes on a timescale set by $\bar W(t)$'s smallest singular value,
substantially faster than $\bar W$'s own evolution. Uniform invertibility of
$\bar W\Sxx\bar W^\top$ (Corollary~\ref{cor:pd-lower}) lets us substitute the
quasi-static decoder $V \leftarrow V_{\mathrm{qs}}(\bar W)$, eliminating $V$
analytically and reducing the joint $(\bar W, V)$ dynamics to an autonomous ODE in
$\bar W$ alone (Proposition~\ref{prop:quasistatic}).

Projecting this ODE onto the generalised eigenbasis $\{v_r^*\}_{r=1}^d$ of
$(\Syx, \Sxx)$, under aligned balanced initialisation $\bar W(t)$ remains diagonal up
to a small off-diagonal leakage $c_{rs}(t)$, uniformly bounded by a
fundamental-theorem-of-calculus and Cauchy--Schwarz argument (Lemma~\ref{lem:offdiag}),
yielding a system of decoupled scalar Bernoulli ODEs (Proposition~\ref{prop:diagonal-ode}).
Diagonalising a Bernoulli-type ODE this way is classical, and we do not claim
otherwise --- it is the one technical step the whole paper leans on without
re-deriving, the step at which simultaneous-diagonalisability of $\Sigma^{xx}$ and
$\Sigma^{yx}$ would otherwise be needed. What is new is what the resulting closed form
is used for: turning a purely qualitative dynamics fact (\emph{which} feature is
learned first) into an explicit, rate-optimal, signed estimator computable from the
trajectory alone.

Each scalar ODE exhibits three qualitatively distinct regimes dictated by
$\rhostar$'s sign (Theorem~\ref{thm:trichotomy}), whose asymptotic shape uniquely
identifies $\operatorname{sign}(\rhostar)$ (Corollary~\ref{cor:sign-id}). On the
positive branch, the plateau value yields $\rhostar$ directly
(Proposition~\ref{prop:plateau}); on the negative branch, late-time decay yields a
projected-covariance estimate (Proposition~\ref{prop:neg-lambda}). The two branches
combine into the unified estimator (Theorem~\ref{thm:headline-formal}).

In practice only $n$ samples are available: replacing the population covariances with
sample estimates $(\hat\Syx,\hat\Sxx)$ introduces basis perturbation, bounded via
matrix Bernstein concentration plus a Wedin-type generalized-eigenvector perturbation
bound (Corollary~\ref{cor:finite-sample-end}). The chain is then packaged as an
algorithm and validated empirically (\S\S\ref{sec:algorithm}--\ref{sec:experiments}).

\subsection*{Outline}
\label{ssec:outline}

\S\S\ref{sec:quasistatic}--\ref{sec:experiments} carry out the
four-step argument sketched above, in the same order.
\S\ref{sec:related} situates the result against prior work;
\S\ref{sec:discussion} discusses the linear-to-non-linear transition
and open questions. Appendix~\ref{app:lean} catalogues the Lean
signatures and the two named axioms; Appendix~\ref{app:finite-sample-derivation}
gives proofs deferred from the main text, including the full
finite-sample perturbation derivation.
